\chapter{Technische Realisierung der Testumgebung}
\label{chap:implementierung}

Dieses Kapitel dokumentiert die technische Implementierung der \emph{BuyIn Todo Sandbox}, die als experimentelle Plattform für diese Studie dient. Es wird erläutert, wie die Brownfield-Umgebung architektonisch aufgebaut ist, um realistische Unternehmensanforderungen zu simulieren, und welche technischen Maßnahmen ergriffen wurden, um eine reproduzierbare Datenerhebung zu gewährleisten.

\section{Anforderungen an die Sandbox}

Um valide Aussagen über Vibe-Coding in Unternehmenskontexten treffen zu können, musste die Testumgebung folgende Anforderungen erfüllen:

\begin{enumerate}
    \item \textbf{Realismus (Brownfield):} Die Codebasis darf nicht leer sein (\enquote{Greenfield}), sondern muss bestehende Strukturen, technische Schulden und Konventionen enthalten, die der Agent und der Mensch verstehen müssen.
    \item \textbf{Komplexität:} Der Tech-Stack muss modern, aber komplex genug sein, um triviale Lösungen zu verhindern.
    \item \textbf{Isolierbarkeit:} Die Umgebung muss für jeden Teilnehmenden identisch zurückgesetzt werden können (Containerisierung).
    \item \textbf{Messbarkeit:} Das System muss automatisierte Tests bereitstellen, um den funktionalen Fortschritt objektiv zu messen.
\end{enumerate}

\section{Architektur und Tech-Stack: Kognitive Dimensionen}

Die Sandbox ist als klassische Full-Stack-Webanwendung konzipiert. Die Architekturwahl hat direkte Auswirkungen auf die kognitive Belastung der Teilnehmenden, analysiert durch das Framework der \emph{Cognitive Dimensions of Notation} \parencite{green1989cognitive}.

\subsection{Backend: Node.js \& Express (Layered Architecture)}
Das Backend basiert auf Node.js mit dem Express-Framework und folgt einer strikten Schichtenarchitektur:
\begin{itemize}
    \item \textbf{Routes:} Definition der API-Endpunkte.
    \item \textbf{Controllers:} Verarbeitung der HTTP-Requests und Validierung.
    \item \textbf{Services:} Kapselung der Geschäftslogik.
    \item \textbf{Repositories:} Datenzugriffsschicht.
\end{itemize}
\textbf{Kognitive Dimension - Viscosity (Zähigkeit):} Diese Trennung erhöht die \emph{Viscosity} für Novizen, da eine Änderung an mehreren Stellen (Datei-Hopping) durchgeführt werden muss. Für Copilot hingegen ist diese Struktur vorteilhaft, da sie klare semantische Grenzen zieht, die das Modell über den Dateinamen und Pfadkontext erkennen kann.

\subsection{Frontend: React \& Vite}
Das Frontend nutzt React mit TypeScript. Die Komponenten sind funktional aufgebaut und nutzen Hooks.
\textbf{Kognitive Dimension - Hidden Dependencies:} Die Verwendung von \emph{Custom Hooks} und \emph{Context API} erzeugt \emph{Hidden Dependencies}. Ein Low-Coder sieht im UI-Code nicht sofort, woher die Daten kommen. Dies testet die Fähigkeit des Agenten, diese Abhängigkeiten aufzulösen, und die Fähigkeit des Nutzers, dem Agenten den richtigen Kontext zu geben.

\section{Technische Funktionsweise von Copilot in der Sandbox}

Um zu verstehen, warum die Sandbox Copilot fordert, muss die technische Funktionsweise betrachtet werden.

\subsection{Prompt Construction und Context Retrieval}
Copilot sendet nicht nur den Code vor dem Cursor an das Modell. Der \emph{Copilot Proxy} in der IDE sammelt Kontext aus:
\begin{itemize}
    \item \textbf{Jaccard Similarity:} Offene Tabs werden nach Ähnlichkeit zum aktuellen Code gescannt.
    \item \textbf{Semantic Search:} Vektorisierte Suche im Projekt (bei neueren Versionen).
\end{itemize}
In der Sandbox sind die Aufgaben so gestaltet, dass die Lösung oft in einer \emph{anderen} Datei liegt (z.B. muss der Controller angepasst werden, um den Service zu nutzen). Wenn der Nutzer diese Datei nicht öffnet, fehlt dem Modell der Kontext. Dies misst indirekt die Kompetenz des Nutzers, dem Agenten die richtigen „Hinweise“ zu geben.

\subsection{Token Prediction und Ranking}
Das Modell generiert mehrere Vorschläge (Beams) und rankt diese nach Wahrscheinlichkeit. In einer Brownfield-Umgebung ist die Wahrscheinlichkeit für „korrekten“ Code höher, wenn er dem Stil der Umgebung entspricht. Ein High-Coder erkennt, ob der Vorschlag den Stil bricht (z.B. \texttt{var} statt \texttt{const}), ein Low-Coder nicht.

\section{Design der Aufgaben (Workload Grading)}

Die Aufgaben sind so gestaffelt, dass sie die kognitive Last sukzessive erhöhen:

\begin{table}[h]
    \centering
    \begin{tabular}{|l|l|l|}
    \hline
    \textbf{Task} & \textbf{Beschreibung} & \textbf{Kognitive Anforderung} \\ \hline
    1 & Persistenz (JSON) & Verstehen des Repository-Patterns (Backend). \\ \hline
    2 & User-Accounts & Einführung neuer Entitäten, Anpassung mehrerer Layer. \\ \hline
    3 & Kategorien (UI) & Frontend-Logik, State-Management. \\ \hline
    4 & Attachments & Umgang mit Binärdaten, API-Erweiterung. \\ \hline
    5 & Kalender & Komplexe Logik, Datumsberechnung (hoher Intrinsic Load). \\ \hline
    6 & Email-Verifikation & Integration externer Konzepte (Mocking). \\ \hline
    7 & Performance & Optimierung, kein funktionales Feature (Experten-Task). \\ \hline
    8 & ICS-Export & Standard-Format, reines Mapping (Erholung). \\ \hline
    \end{tabular}
    \caption{Aufgaben-Design und kognitive Anforderungen}
    \label{tab:task-design}
\end{table}

\section{Automatisierung der Evaluation}

Um die Auswertung der Ergebnisse zu standardisieren, wurde ein Evaluations-Framework entwickelt.

\subsection{Das Analyse-Skript (\texttt{run.js})}
Ein Node.js-Skript automatisiert die Prüfung der Lösungen. Es führt folgende Schritte aus:
\begin{enumerate}
    \item \textbf{Test-Execution:} Ausführung der Unit-Tests (Vitest) und End-to-End-Tests (Playwright).
    \item \textbf{Linting:} Prüfung auf Code-Style-Verletzungen mittels ESLint.
    \item \textbf{Report-Generierung:} Zusammenfassung der Ergebnisse in einer JSON-Datei.
\end{enumerate}

\subsection{Grenzen der LLM-gestützten Code-Analyse}
Zusätzlich zur rein funktionalen Prüfung wird ein LLM-Prompt-Template verwendet, um qualitative Aspekte des Codes zu bewerten. Hierbei muss kritisch angemerkt werden, dass LLMs zu \emph{Halluzinationen} neigen können, also Fehler finden, die nicht existieren, oder echte Probleme übersehen. Daher dient dieser Score nur als Indikator und wird stichprobenartig manuell validiert.

\section{Relevanz für den Unternehmenskontext (BuyIn)}

Die gewählte Architektur und Methodik spiegelt reale Herausforderungen bei BuyIn wider:
\begin{itemize}
    \item \textbf{Legacy-Integration:} Wie bei BuyIn müssen Entwickler oft in gewachsenen Strukturen arbeiten, die nicht perfekt dokumentiert sind. Vibe-Coding kann hier helfen, den „impliziten Kontext“ schneller zu erschließen.
    \item \textbf{Security \& Compliance:} Der Einsatz von Copilot in Enterprise-Umgebungen birgt Risiken (Datenabfluss, Lizenzverletzungen). Die Sandbox simuliert dies durch Aufgaben, bei denen sensible Daten (User-Accounts) behandelt werden müssen.
    \item \textbf{Skalierbarkeit:} Die modulare Struktur der Sandbox zeigt, wie Copilot in größeren Teams (wo Modularisierung Pflicht ist) performt, im Gegensatz zu kleinen Skript-Projekten.
\end{itemize}

\section{Containerisierung und Reproduzierbarkeit}

Die gesamte Umgebung ist mittels Docker containerisiert. Dies stellt sicher, dass alle Teilnehmenden – unabhängig von ihrem Betriebssystem – unter identischen Bedingungen arbeiten und technische Störfaktoren (z.B. unterschiedliche Node-Versionen) eliminiert werden.
